\chapter{正則関数}

\section{多変数正則関数}

\subsection*{記法と準備}

$n$ 次元複素空間を
\[
  \C^n = \{\, z = (z_1, \dots, z_n) \mid z_\nu \in \C \,\}
\]
とし，各成分を $z_\nu = x_\nu + i y_\nu$（$x_\nu, y_\nu \in \R$，$i = \sqrt{-1}$）と
書く．実部・虚部をとる対応 $z \mapsto (x_1, y_1, \dots, x_n, y_n)$ により
$\C^n \cong \R^{2n}$ とみなす．

$z = (z_1, \dots, z_n)$ のノルムを
\[
  |z| = \Bigl( \sum_{\nu=1}^{n} |z_\nu|^2 \Bigr)^{1/2}
\]
で定める．$z, w \in \C^n$ に対して次が成り立つ：
\begin{align}
  |z_\nu| &\le |z| \le \sum_{\nu=1}^n |z_\nu|, \\
  |z + w| &\le |z| + |w|, \\
  \bigl|\, |z| - |w| \,\bigr| &\le |z - w|.
\end{align}


\subsection*{正則関数}

\begin{definition}\label{def:holomorphic}
$D \subseteq \C^n$ とする．$f(z) = f(z_1, \dots, z_n)$ が $D$ において
\emph{正則関数}（holomorphic function）であるとは，$f$ が $D$ 上連続で，かつ
各変数 $z_\nu$ について（他の変数を固定するごとに）正則であることをいう．
\end{definition}

\subsection*{領域と多重円板}

\begin{definition}[領域]
$\C^n$ の連結開集合を\emph{領域}（domain）という．
\end{definition}

\begin{note}
位相空間 $X$ が\emph{連結}であるとは，$X = U \cup V$，$U \cap V = \varnothing$，
$U, V$ は開集合，という分割が $U = \varnothing$ または $V = \varnothing$ の場合に
限ること（すなわち空でない二つの開集合に分割できないこと）をいう．$\C^n$ の
開集合については，連結であることと弧状連結（任意の二点が領域内の道で結べること）
であることは同値である．
\end{note}

\begin{definition}[多重円板]
$a = (a_1, \dots, a_n) \in \C^n$ $r = (r_1, \dots, r_n)$
（$r_\nu > 0$）に対し，
\[
  P(a, r) = \{\, z \in \C^n \mid |z_\nu - a_\nu| < r_\nu,\ \nu = 1, \dots, n \,\}
\]
を中心 $a$，半径 $r$ の\emph{多重円板}（polydisc）という．
\end{definition}

\begin{note}
最大値ノルム $\|z\| = \max_{1 \le \nu \le n} |z_\nu|$ を用いると，多重円板は
$\|z - a\|$ に関する「球」$\{\, z \mid \|z - a\| < r \,\}$（$r_1 = \dots = r_n = r$
のとき）にあたる．
\end{note}

\subsection*{Wirtinger 微分(複素微分)と Cauchy--Riemann 方程式}

$z_\nu = x_\nu + i y_\nu$ に対し，\emph{Wirtinger 微分}を
\[
  \frac{\partial}{\partial z_\nu}
    = \frac{1}{2}\Bigl( \frac{\partial}{\partial x_\nu} - i \frac{\partial}{\partial y_\nu} \Bigr),
  \qquad
  \frac{\partial}{\partial \bar z_\nu}
    = \frac{1}{2}\Bigl( \frac{\partial}{\partial x_\nu} + i \frac{\partial}{\partial y_\nu} \Bigr)
\]
で定める．以下，$C^1$ 級関数 $f$ を実部・虚部に分けて $f = u + i v$
（$u, v$ は実数値関数）と書く．

\begin{definition}[正則, Cauchy--Riemann 方程式]
$f = u + i v$ が変数 $z_\nu$ について\emph{正則}であるとは，
\[
  \frac{\partial f}{\partial \bar z_\nu} = 0
\]
が成り立つことをいう．これは Cauchy--Riemann 方程式
\[
  \frac{\partial u}{\partial x_\nu} = \frac{\partial v}{\partial y_\nu}, \qquad
  \frac{\partial u}{\partial y_\nu} = -\frac{\partial v}{\partial x_\nu}
\]
と同値である．
\end{definition}

\begin{proof}
定義より
\[
  \frac{\partial f}{\partial \bar z_\nu}
  = \frac{1}{2}\Bigl( \frac{\partial}{\partial x_\nu} + i \frac{\partial}{\partial y_\nu} \Bigr)(u + i v)
  = \frac{1}{2}\Bigl[ \Bigl( \frac{\partial u}{\partial x_\nu} - \frac{\partial v}{\partial y_\nu} \Bigr)
      + i \Bigl( \frac{\partial v}{\partial x_\nu} + \frac{\partial u}{\partial y_\nu} \Bigr) \Bigr].
\]
複素数が $0$ であることは実部と虚部がともに $0$ であることと同値だから，
$\partial f / \partial \bar z_\nu = 0$ は上の Cauchy--Riemann 方程式と同値である．
\end{proof}

\subsection*{べき級数の収束}

\begin{note}
$f(z)$ が $W$ の各点 $a$ の近傍で収束べき級数
\[
  f(z) = \sum_{k_1, \dots, k_n \ge 0}
    c_{k_1 \cdots k_n} (z_1 - a_1)^{k_1} \cdots (z_n - a_n)^{k_n}
\]
に展開できるとき，$f(z)$ は正則関数である．
\end{note}

\begin{proposition}\label{cl:abel}
べき級数 $\sum_{k_1, \dots, k_n \ge 0} c_{k_1 \cdots k_n}
(z_1 - a_1)^{k_1} \cdots (z_n - a_n)^{k_n}$ がある点 $w$ で収束すれば，これは
\[
  |z_\nu - a_\nu| < |w_\nu - a_\nu| \qquad (\nu = 1, \dots, n)
\]
を満たす $z$ において絶対収束する．
\end{proposition}

\begin{proof}
$a = 0$ としてよい（$z_\nu - a_\nu$ を改めて $z_\nu$ とおけばよい）．級数が $w$ で
収束するから，その各項は有界であり，
\[
  |c_{k_1 \cdots k_n} \, w_1^{k_1} \cdots w_n^{k_n}| \le C
  \qquad (\forall\, k_1, \dots, k_n)
\]
となる定数 $C < \infty$ がとれる．$|z_\nu| < |w_\nu|$ のとき
\[
  |c_{k_1 \cdots k_n} \, z_1^{k_1} \cdots z_n^{k_n}|
  = |c_{k_1 \cdots k_n} \, w_1^{k_1} \cdots w_n^{k_n}|
    \prod_{\nu=1}^n \left| \frac{z_\nu}{w_\nu} \right|^{k_\nu}
  \le C \prod_{\nu=1}^n \left| \frac{z_\nu}{w_\nu} \right|^{k_\nu}
\]
だから，
\[
  \sum_{k_1, \dots, k_n \ge 0} |c_{k_1 \cdots k_n} \, z_1^{k_1} \cdots z_n^{k_n}|
  \le C \prod_{\nu=1}^n \sum_{k=0}^\infty \left| \frac{z_\nu}{w_\nu} \right|^{k}
  = C \prod_{\nu=1}^n \frac{1}{1 - |z_\nu / w_\nu|} < \infty .
\]
よって級数は絶対収束する．
\end{proof}

\subsection*{Theorem 1：べき級数展開}

\begin{theorem}\label{thm:power-series}
$W \subseteq \C^n$ を領域とする．$f(z)$ が $W$ で連続で，かつ各変数 $z_\nu$ について
（他の変数を固定するごとに）正則ならば，$f(z)$ は $n$ 変数 $z = (z_1, \dots, z_n)$ の
正則関数である．
\end{theorem}

\begin{proof}
$a$ を $W$ の任意の点とする．$W$ は開集合だから，
\[
  \{\, z \mid |z_\nu - a_\nu| \le r_\nu,\ \nu = 1, \dots, n \,\} \subset W
\]
となる $r = (r_1, \dots, r_n)$ がとれる．$f$ は $z_1$ について正則だから，$z_1$ 平面で
Cauchy の積分公式を用いて
\[
  f(z_1, z_2, \dots, z_n)
  = \frac{1}{2\pi i} \oint_{|w_1 - a_1| = r_1}
      \frac{f(w_1, z_2, \dots, z_n)}{w_1 - z_1}\, dw_1 .
\]
各変数についてこれを繰り返す．$f$ は連続だから反復積分は重積分に等しく（積分の
順序は交換できる），
\[
  f(z) = \frac{1}{(2\pi i)^n}
    \oint_{|w_1 - a_1| = r_1} \!\!\!\! \cdots \oint_{|w_n - a_n| = r_n}
    \frac{f(w_1, \dots, w_n)}{(w_1 - z_1) \cdots (w_n - z_n)}\, dw_1 \cdots dw_n .
\]
ここで $|z_\nu - a_\nu| < r_\nu = |w_\nu - a_\nu|$ に注意すると，各因子は
\[
  \frac{1}{w_\nu - z_\nu}
  = \frac{1}{w_\nu - a_\nu} \cdot
    \frac{1}{1 - \dfrac{z_\nu - a_\nu}{w_\nu - a_\nu}}
  = \sum_{k=0}^\infty \frac{(z_\nu - a_\nu)^k}{(w_\nu - a_\nu)^{k+1}}
\]
と展開できる．これらを代入して項別に積分すると，
\[
  f(z) = \sum_{k_1, \dots, k_n \ge 0}
    c_{k_1 \cdots k_n} (z_1 - a_1)^{k_1} \cdots (z_n - a_n)^{k_n} ,
\]
ただし
\[
  c_{k_1 \cdots k_n} = \frac{1}{(2\pi i)^n}
    \oint_{|w_1 - a_1| = r_1} \!\!\!\! \cdots \oint_{|w_n - a_n| = r_n}
    \frac{f(w_1, \dots, w_n)}
         {(w_1 - a_1)^{k_1 + 1} \cdots (w_n - a_n)^{k_n + 1}}\, dw_1 \cdots dw_n
\]
である．よって $f$ は $W$ の各点の近傍で収束べき級数に展開され，$n$ 変数の
正則関数である．
\end{proof}
